% Theorem environments, numbered together (Theorem 1, Lemma 2, ...).
% \declaretheorem (thmtools) keeps cleveref's names right for shared counters:
% \cref{def:x} -> "Definition 1", not "Theorem 1".
\declaretheorem[style=plain]{theorem}
\declaretheorem[style=plain, sibling=theorem]{lemma}
\declaretheorem[style=plain, sibling=theorem]{proposition}
\declaretheorem[style=definition, sibling=theorem]{definition}
\declaretheorem[style=definition]{assumption}
\declaretheorem[style=remark, sibling=theorem]{remark}

% Sets and operators.
\newcommand{\R}{\mathbb{R}}
\newcommand{\N}{\mathbb{N}}
\newcommand{\E}{\mathbb{E}}
\DeclareMathOperator*{\argmin}{arg\,min}
\DeclareMathOperator*{\argmax}{arg\,max}
\DeclarePairedDelimiter{\abs}{\lvert}{\rvert}
\DeclarePairedDelimiter{\norm}{\lVert}{\rVert}

% Vectors and matrices: bold, e.g. \vect{x}, \mat{A}.
\newcommand{\vect}[1]{\bm{#1}}
\newcommand{\mat}[1]{\bm{#1}}

% Inline note for unfinished parts; remove all before sharing.
\newcommand{\todo}[1]{\textcolor{red}{[TODO: #1]}}
